\Needspace{12\baselineskip}
\begin{proposition}\label{prop:smooth-obstacle-support}\index{Contact set!exponential approximation}
Let $\omega$ be Hermitian and $f\in C^\infty(X,\mathbb R)$. For $\beta>0$, denote by $u_\beta$ the smooth strictly $\omega$-plurisubharmonic solution of
\[
 \omega_{u_\beta}^n=\mathrm{e}^{\beta(u_\beta-f)}\omega^n.
\]
These solutions converge uniformly as $\beta\to+\infty$ to $p=P_\omega(f)$. Moreover,
\[
 \mathbf1_{\{p<f\}}\omega_p^n=0.
\]
\end{proposition}
\begin{proof}
Existence is \cref{thm:smooth-proper-existence}. Choose $A\geq0$ with $\ddc f\leq A\omega$. At a maximum of $u_\beta-f$,
\[
 \omega_{u_\beta}\leq\omega+\ddc f\leq(1+A)\omega.
\]
\Needspace{5\baselineskip}
The equation gives
\[
 u_\beta\leq f+\frac{C_f}{\beta},\qquad C_f=n\log(1+A).
\]
Thus $u_\beta-C_f/\beta$ is an envelope competitor, and
$u_\beta\leq p+C_f/\beta$.

The function $p$ is continuous by \cref{prop:continuous-obstacle-envelope}. Choose smooth strictly $\omega$-plurisubharmonic $v_j\downarrow p$ by \cref{thm:bounded-exact-regularization}; convergence is uniform by compactness. For any $\delta>0$, choose $j$ so that
\[
 p\leq v_j\leq p+\delta/2.
\]
Then $v_j-\delta\leq f$. Positivity and compactness give a number $c_j>0$ such that $\omega_{v_j}^n\geq c_j\omega^n$. Put $C_j=\max(0,-\log c_j)$. The function
\[
 w_{j,\beta}=v_j-\delta-C_j/\beta
\]
satisfies
\[
 \omega_{w_{j,\beta}}^n=\omega_{v_j}^n
 \geq \mathrm{e}^{-C_j}\omega^n
 \geq \mathrm{e}^{\beta(w_{j,\beta}-f)}\omega^n.
\]
Smooth exponential comparison gives $w_{j,\beta}\leq u_\beta$. Consequently
\[
 p-\delta-C_j/\beta\leq u_\beta\leq p+C_f/\beta.
\]
First let $\beta$ tend to infinity with $j$ fixed, so no uniform lower bound on $c_j$ is needed, and then let $\delta\downarrow0$. This proves uniform convergence. Bounded Bedford--Taylor continuity gives weak convergence of the measures.

On a compact subset of the open set $\{p<f\}$, uniform convergence gives $u_\beta-f\leq-\epsilon$ for all large $\beta$, with some $\epsilon>0$. There the measures have densities at most $\mathrm{e}^{-\beta\epsilon}$. Testing with continuous functions of compact support in that open set and passing to the bounded-product limit gives $\omega_p^n=0$ there.
\end{proof}

\begin{proposition}
\label{prop:lsc-obstacle-support}
If $\omega$ is Hermitian and $f$ is bounded and lower semicontinuous, then $P_\omega(f)$ is bounded, lies below $f$ quasi-everywhere, and its measure is concentrated on $\{P_\omega(f)=f\}$.
\end{proposition}
\begin{proof}
Suppose first that $f$ is continuous. Choose smooth $f_j\to f$ uniformly. The envelope Lipschitz bound gives $p_j=P_\omega(f_j)\to p=P_\omega(f)$ uniformly. On every compact subset of $\{p<f\}$, the inequalities $p_j<f_j$ hold on a neighborhood for all large $j$. The preceding proposition and bounded-product convergence then give $\omega_p^n=0$ on $\{p<f\}$.

For lower semicontinuous $f$, choose continuous $f_j\uparrow f$, with common upper and lower bounds. For example, for a fixed compatible distance $d$ one may take
\[
 f_j(x)=\inf_{y\in X}\bigl(f(y)+j d(x,y)\bigr).
\]
These functions are Lipschitz and increase to $f$: if convergence failed at $x$, approximate minimizers would approach $x$, contradicting lower semicontinuity.

The functions $p_j=P_\omega(f_j)$ increase, stay bounded, and have a bounded $\omega$-plurisubharmonic regularized limit $p_\infty=(\lim_jp_j)^*$. A countable union of pluripolar exceptional sets shows that $p_\infty\leq f$ quasi-everywhere. Thus $p_\infty\leq P_\omega(f)$.

For fixed $k$, the set $\{p_\infty<f_k\}$ is open. If $j\geq k$, then $p_j\leq p_\infty<f_k\leq f_j$ there, so the continuous-obstacle case gives zero measure on that open set. Increasing bounded Bedford--Taylor convergence applies after adding a fixed smooth local correction whose Hessian dominates $\omega$ and expanding against the resulting smooth background. Together with the weak-limit inequality for open sets, it yields
\[
 \omega_{p_\infty}^n(\{p_\infty<f_k\})=0.
\]
Taking the union over $k$ gives zero measure on $\{p_\infty<f\}$, hence also on $\{p_\infty<P_\omega(f)\}$ outside a pluripolar set. By bounded domination, \cref{thm:bounded-strict-domination}, $p_\infty\geq P_\omega(f)$. They are therefore equal, and the assertion follows.
\end{proof}

\begin{theorem}\label{thm:bounded-contact-support}
Let $\theta\geq0$ be smooth and let $f$ be bounded and quasi-continuous. Then $p=P_\theta(f)$ is bounded, $p\leq f$ quasi-everywhere, and
\[
 \mathbf1_{\{p<f\}}\mathrm{MA}_\theta(p)=0.
\]
\end{theorem}
\begin{proof}
First suppose $\theta=\omega$ is Hermitian. Choose $f_j\downarrow f$ quasi-everywhere as in \cref{lem:quasicont-lsc-approximation}. The bounded functions $p_j=P_\omega(f_j)$ decrease to $p$. Indeed their limit is a competitor for $p$, since it lies below $f$ quasi-everywhere; the reverse inequality follows from monotonicity.

For fixed $k$, put $H_k=\min\{1,(f-p_k)_+\}$. It is bounded and quasi-continuous. Where $H_k>0$ and $j\geq k$, we have $p_j\leq p_k<f\leq f_j$, outside a pluripolar set. The lower semicontinuous case therefore gives
\[
 \int_X H_k\,\omega_{p_j}^n=0.
\]
The weighted bounded-product convergence in \cref{sec:hermitian-capacity} gives
$\int_XH_k\,\omega_p^n=0$. The sets $\{p_k<f\}$ increase to $\{p<f\}$, proving support.

For general semipositive $\theta$, set $\theta_j=\theta+j^{-1}\omega$ with $\omega$ Hermitian. The functions $p_j=P_{\theta_j}(f)$ have the common bounds $\inf_Xf\leq p_j\leq\sup_Xf$ and decrease. Their limit is $\theta$-plurisubharmonic by distributional convergence, lies below $f$ quasi-everywhere, and hence equals $p=P_\theta(f)$. For each fixed $k$, the same coefficient $H_k=\min\{1,(f-p_k)_+\}$ gives
\[
 \int_XH_k\,\mathrm{MA}_{\theta_j}(p_j)=0\qquad(j\geq k).
\]
All $p_j$ are bounded potentials for the fixed Hermitian form $\theta+\omega$. Expand the products with $\theta_j$ against that fixed background. Bounded-product convergence and uniform capacity control pass the displayed identity to $\mathrm{MA}_\theta(p)$, including its quasi-continuous coefficient. Exhausting $\{p<f\}$ as before completes the proof.
\end{proof}
